\documentclass[14pt]{extarticle}
\usepackage{preamble}
\geometry{letterpaper, landscape, margin=0.5in}
\usepackage{qrcode}
\renewcommand{\answer}[1]{}
\setlength{\parskip}{4pt}

\begin{document}
\pagestyle{empty}

\daybanner{Unit 1 \textperiodcentered\ Topic 1.5 \textperiodcentered\ TODAY'S PROBLEM}{Same Data, Different Bins}

\noindent\begin{minipage}[t]{0.57\linewidth}
\vspace{0pt}
Suppose a student recorded steps for 30 days. Histograms F and C show the same data with 1-thousand- and 5-thousand-step bins. A coach asks, ``Does this student move differently on school days and weekends?''\par\smallskip \textbf{Priya:} ``Use C. F has too many short and empty bars; C gives the clearer summary.''\par \textbf{Leo:} ``Use F. C combines values that should not be treated as one group.''\par
\end{minipage}\hfill
\begin{minipage}[t]{0.40\linewidth}
\vspace{0pt}
\begin{center}
\noindent
\begin{minipage}[t]{0.46\linewidth}
\centering
\textbf{F: 1-thousand-step bins}\par\smallskip
\begin{tikzpicture}
\begin{axis}[
  width=\linewidth, height=1.40in,
  ybar interval, bar shift=0pt,
  xmin=4, xmax=20, ymin=0, ymax=8,
  xtick={4,6,8,10,12,14,16,18,20}, ytick={0,2,4,6,8},
  xlabel={Daily steps (thousands)}, ylabel={Days},
  tick label style={font=\scriptsize}, label style={font=\scriptsize},
  ymajorgrids, x tick label as interval=false
]
\addplot[fill=calloutblue, draw=dayblue] coordinates {
  (4,0) (5,5) (6,6) (7,7) (8,3) (9,0) (10,0) (11,3)
  (12,3) (13,2) (14,0) (15,0) (16,0) (17,0) (18,0) (19,1) (20,0)
};
\end{axis}
\end{tikzpicture}
\end{minipage}\hfill
\begin{minipage}[t]{0.46\linewidth}
\centering
\textbf{C: 5-thousand-step bins}\par\smallskip
\begin{tikzpicture}
\begin{axis}[
  width=\linewidth, height=1.40in,
  ybar interval, bar shift=0pt,
  xmin=0, xmax=20, ymin=0, ymax=23,
  xtick={0,5,10,15,20}, ytick={0,5,10,15,20},
  xlabel={Daily steps (thousands)}, ylabel={Days},
  tick label style={font=\scriptsize}, label style={font=\scriptsize},
  ymajorgrids, x tick label as interval=false
]
\addplot[fill=calloutblue, draw=dayblue] coordinates {
  (0,0) (5,21) (10,8) (15,1) (20,0)
};
\end{axis}
\end{tikzpicture}
\end{minipage}
\end{center}
\end{minipage}

\begin{firsttakebox}{FIRST TAKE (5 min, on your sheet)}
Which histogram is more useful to the coach, F or C? Name one visible feature.
\end{firsttakebox}

\begin{description}[leftmargin=1.15in, labelwidth=1in, itemsep=2pt, topsep=2pt]
  \item[\dokbadge{1}~(a)] State each bin width and the number of days in F's 19--20-thousand interval.
  \item[\dokbadge{2}~(b)] Describe F in context, including clusters, gaps, and the unusual value, with units.
  \item[\dokbadge{3}~(c)] Adjudicate Priya's and Leo's recommendations. Choose the better first look for the coach, cite a feature that vanishes in the other display, and explain what its reader could wrongly conclude. State what information is still needed before naming either cluster school days or weekends. Give one specific question for which C is better.
\end{description}

\vfill
\noindent\begin{minipage}{0.84\linewidth}
\textbf{Video + follow-along:} \href{https://robjohncolson.github.io/apstats-live-worksheet/u1_lesson5_live.html}{\texttt{\detokenize{u1_lesson5_live.html}}} (scan the code)\par
\textbf{Finish (a)--(c) after the video. Turn in your sheet.}
\end{minipage}\hfill
\qrcode[height=0.75in]{https://robjohncolson.github.io/apstats-live-worksheet/u1_lesson5_live.html}

\end{document}
